\scanpage{44}
b) $v_1$ et $v_2$ représentent, par abus de langage, les actions de $v_1\in V_1$ et de $v_2\in V_2$ sur $E$.

c) $v_1\cos\theta+\varepsilon^i\sin\theta$, $i=1,2$, $\theta\in[0,\pi/2]$, représentent, encore par abus de langage, les graduations $v_1\cos+\varepsilon^i\sin\theta$ au point de $S^+(V_1+1)$ de coordonnées $v_1\cos\theta+\varepsilon\sin\theta$, $v_1$ appartenant à $S(V_1)$ et $\varepsilon$ représentant le vecteur unité de $R$.

\uuline{Théorème 1.1} $^{(2)}$: L'homomorphisme
\[
t:K^{V_1+V_2}(X)\longrightarrow K^{V_2}(B(V_1),S(V_1))
\]
est un isomorphisme.

1.4\quad Soit $\operatorname{Spin}(p,q)$ le sous-groupe du groupe multiplicatif $(C^{p,q})^*$ formé des éléments $x$ de degré zéro tels que
\[
\begin{aligned}
\text{a)}\quad&\forall\ y\in R^{p+q},\quad xyx^{-1}\in R^{p+q}\\
\text{b)}\quad&x^*x=1,
\end{aligned}
\]
$*$ désignant l'antiinvolution de $C^{p,q}$ définie sur les générateurs $e_i$, $i=1,\ldots,p+q$, par $e_i^*=-e_i$.
Pour un élément $x$ de $\operatorname{Spin}(p,q)$ la transformation $\sigma(x)$ définie par $\sigma(x)y=xyx^{-1}$ appartient en fait à $\operatorname{SO}(p,q)$ et on a une suite exacte
\diagram{\begin{tikzcd}[ampersand replacement=\&]
0\arrow[r]\&Z_2\arrow[r]\&\operatorname{Spin}(p,q)\arrow[r,"\sigma"]\&\operatorname{SO}(p,q)\arrow[r]\&0
\end{tikzcd}}
Pour plus d'informations sur le groupe $\operatorname{Spin}(p,q)$, voir $^{(2)}$.

Soit maintenant $V$ un fibré vectoriel réel muni d'une forme quadratique non dégénérée de type $(q,p)$, i.e. dont le groupe structural est réduit à $\operatorname{O}(p,q)$. On dira que le fibré $V$ est spinoriel de type $\operatorname{Spin}(p,q)$ si on se donne un fibré principal $P$ de groupe $\operatorname{Spin}(p,q)$ tel que $V\approx P\,x_{\operatorname{Spin}(p,q)}R^{p+q}$, $\operatorname{Spin}(p,q)$ opérant sur $R^{p+q}$ par la représentation $\sigma$. Si $\mathcal C$ est une catégorie de Banach arbitraire, on définit maintenant une équivalence de catégories
\mspage{43}
